\section{Related Work and Positioning}
\label{sec:related}
Time-distributed FFT execution has a substantial history. Lo and Lee construct
butterfly work during acquisition and also combine scheduling with recursive
reuse across shifted frames~\cite{lo1998,lo1999,lomoving1999}. In convolution,
Bor\ss{} distributes partition work and staggers filter schedules to control
combined load~\cite{borss2009}, while Hurchalla distributes transform work
across incoming blocks on one thread~\cite{hurchalla2010}. Battenberg and
Avi\v{z}ienis schedule forward transforms, spectral products, and inverse
transforms across callbacks, using FFTW for leaf transforms and comparing
cooperative with preemptive execution~\cite{battenberg2011}. Thus both complete
processing chains and optimized transform kernels have precedents within
cooperative scheduling. Wefers treats manual preemption and transform
subdivision more broadly~\cite{wefers2015}; suspension granularity remains an
implementation choice rather than an inherent advantage of butterfly-level
execution. Liu, Yan, and Zou explicitly study time-distributed FFT and inverse
FFT execution for online control~\cite{liu2017}. Neither transform suspension
nor extending it to an inverse chain is a new principle of this report.

For spectral visualization, Pr\r{u}\v{s}a and Holighaus use a worker thread
and distinguish computation, polling, and repaint delay~\cite{prusa2016}.
Their filter-bank representation differs from the fixed-window FFT considered
here, but the separation of execution context and visible-output latency is
directly relevant. Sliding, hopping, and feedforward transforms offer another
dimension of choice: reuse across overlapping
windows~\cite{lyons2021,park2014,garrido2016}. Frequency-domain kernels can
support tapered windows in such methods~\cite{rafii2018,eleftheriadis2023}.
Overlap reuse and time distribution can be combined; they are not mutually
exclusive alternatives. Our implementation selects a complete window and
starts a fresh transform for each frame.

Within this literature, the present contribution is an implementation contract
and an empirical study: a complete windowed analyzer with an exact publication
hop, retained input, cancellation, and ownership of published output, evaluated
alongside periodic inverse jobs and complete transform--filter--inverse chains.
The inverse and overlap-save chains are benchmark applications of the same
scheduling principle; they do not introduce a new convolution method.
We measure FFTW, Rack's PFFFT wrapper, and Apple Accelerate/vDSP as practical
optimized batch baselines~\cite{frigo2005,pffft,vdsp}. Matched first-party and
PFFFT hybrid controls distinguish work placement from changing the underlying
FFT implementation. The comparison accounts for input frames, bins, window,
normalization, output processing, numerical acceptance, and layout. It does not
rank the unmeasured worker-thread or overlap-reuse alternatives.
Section~\ref{sec:evaluation} defines the measurement boundaries;
Appendix~\ref{app:related-work} gives the detailed literature comparison.
